%%%%%%%%%%%%%%%%%%%%%%%%%%%%%%%%%%%%%%%%%
% Medium Length Professional CV
%
% This is the shared document shell used for every language build.
% It carries NO candidate data itself - the CLI generates a sibling
% file named curriculo_<lang>.tex (e.g. curriculo_en.tex) that defines
% every \newcommand referenced below (\cvname, \resumoprofissional,
% \skillsRows, \experiencias, etc.), then compiles this file next to it.
%%%%%%%%%%%%%%%%%%%%%%%%%%%%%%%%%%%%%%%%%

\documentclass{resume}

% resume.cls already loads and configures geometry (0.75in/0.6in margins) -
% do not \usepackage{geometry} again here, it causes an "Option clash" error.
\usepackage[utf8]{inputenc}
\usepackage{etoolbox}
\usepackage{xparse}

%----------------------------------------------------------------------------------------
% LANGUAGE DATA FILE
% The CLI replaces the line below with e.g. \input{curriculo_en.tex}
% for whichever language it is currently building.
%----------------------------------------------------------------------------------------

% {{INPUT_CURRICULO}}

%----------------------------------------------------------------------------------------

\name{\cvname}
\address{\cvaddress}
\address{\cvphone \\ \cvemail}
\address{\cvlinkedin}

\begin{document}

%----------------------------------------------------------------------------------------
% SUMMARY
% Sourced from: input/contact.json -> summary.<lang>
%----------------------------------------------------------------------------------------

\begin{rSection}{\tituloResumo}
\resumoprofissional
\end{rSection}

%----------------------------------------------------------------------------------------
% TECHNICAL SKILLS
% Sourced from: input/stacks.json (one row per skill group)
% \skillsRows is generated as a sequence of "\label & item, item, ... \\" lines,
% one per entry in stacks.json, using each group's label.<lang> and items.
%----------------------------------------------------------------------------------------

\begin{rSection}{\tituloCompetencias}

\begin{tabular}{ @{} >{\bfseries}l @{\hspace{6ex}} l }
\skillsRows
\end{tabular}

\end{rSection}

%----------------------------------------------------------------------------------------
% EXPERIENCE
% Sourced from: input/experiences.json
% \experiencias is a csv list of {Company|Start - End|Role|Location|bullets}
% groups, one per job, generated in the same order as the JSON array.
% Each job's optional "tag" becomes a bolded suffix on the role line, and its
% optional "highlight" becomes an extra emphasized bullet at the end.
%----------------------------------------------------------------------------------------

\begin{rSection}{\tituloExperiencia}

\NewDocumentCommand{\processExperience}{>{\SplitArgument{4}{|}}m}{%
  \processExperienceAux#1%
}

\NewDocumentCommand{\processExperienceAux}{mmmmm}{%
  \begin{rSubsection}{#1}{#2}{#3}{#4}
  #5
  \end{rSubsection}
}

\expandafter\forcsvlist\expandafter\processExperience\expandafter{\experiencias}

\end{rSection}

%----------------------------------------------------------------------------------------
% EDUCATION
% Sourced from: input/education.json (first entry; extend this section if you
% need to render more than one degree)
%----------------------------------------------------------------------------------------

\begin{rSection}{\tituloFormacao}

{\bf \formacaoInstituicao} \hfill {\em \formacaoData} \\
\formacaoCurso \\
\formacaoDesc

\end{rSection}

%----------------------------------------------------------------------------------------
% LANGUAGES
% Sourced from: input/languages.json (\idiomaRows is one "\label & level \\"
% line per entry, in the same order as the JSON array)
%----------------------------------------------------------------------------------------

\begin{rSection}{\tituloIdiomas}

\begin{tabular}{ @{} >{\bfseries}l @{\hspace{6ex}} l }
\idiomaRows
\end{tabular}

\end{rSection}

\end{document}
